\section{Introduction}
\label{sec:intro}

\subsection{Two questions}

The first question \cite{q1} asks for an asymptotic expansion, as $y\to0^+$, of the inverse of
\begin{equation}
 f_1(x)=x+x^2(1+\log x),\qquad x>0,
 \label{eq:f1}
\end{equation}
on the branch that maps $(0,\infty)$ to itself. Choosing this branch is essential: near a nonzero complex preimage of $0$, the same expression has an analytic inverse germ of the form
\begin{equation}
 c+\frac{y}{c-1}-\frac{3+2\W_{-1}(-e)}{2(c-1)^3}\,y^2+\OO(y^3),\qquad c=e^{\W_{-1}(-e)-1}\approx-0.1178-0.5332\,i,
 \label{eq:complex-germ}
\end{equation}
Here $c$ is a nonzero complex root of $1+x(1+\log x)=0$, where $f_1$ is analytic with a nonzero derivative (Section~\ref{sec:branch}). The real inverse has a first-order Taylor expansion at $0$, but no second-order Taylor expansion there. Its power--logarithmic expansion is
\begin{equation}
\begin{aligned}
 g_1(y)={}&y-y^2(1+L)+y^3\bigl(2L^2+5L+3\bigr)-y^4\Bigl(5L^3+\tfrac{41}{2}L^2+27L+\tfrac{23}{2}\Bigr)\\
 &+y^5\Bigl(14L^4+\tfrac{241}{3}L^3+\tfrac{335}{2}L^2+151L+\tfrac{299}{6}\Bigr)+\OO\bigl(y^6(1+|L|)^5\bigr),
\end{aligned}
 \label{eq:answer1}
\end{equation}
with $L=\log y$: every power of $y$ carries a polynomial in $\log y$, of degree growing with the power, so the scale of the expansion is neither a Taylor nor a Puiseux scale. The question also asks, as a second example, for
\begin{equation}
 g_2(y)=y-y^{\sqrt2}+\sqrt2\,y^{2\sqrt2-1}-\frac{6-\sqrt2}{2}\,y^{3\sqrt2-2}+\Bigl(\frac{17\sqrt2}{3}-4\Bigr)y^{4\sqrt2-3}+\OO\bigl(y^{5\sqrt2-4}\bigr),
 \label{eq:answer2}
\end{equation}
the inverse of $f_2(x)=x+x^{\sqrt2}$: no logarithms this time, but exponents in the semigroup $1+(\sqrt2-1)\N$, which no rational grid contains.

The second question \cite{q2} concerns the forward asymptotics of $(1+x+x^{\sqrt2})^{\sqrt2}$ as $x\to\infty$. Its leading term is $x^2$, but the subsequent terms require ordering two incommensurable positive gaps. The expansion is
\begin{equation}
\begin{aligned}
 (1+x+x^{\sqrt2})^{\sqrt2}={}&x^2+\sqrt2\,x^{3-\sqrt2}+\Bigl(1-\tfrac1{\sqrt2}\Bigr)x^{4-2\sqrt2}+\Bigl(\tfrac{2\sqrt2}{3}-1\Bigr)x^{5-3\sqrt2}\\
 &+\sqrt2\,x^{2-\sqrt2}+\tfrac{13-9\sqrt2}{12}\,x^{6-4\sqrt2}+(2-\sqrt2)\,x^{3-2\sqrt2}+\OO\bigl(x^{7-5\sqrt2}\bigr),
\end{aligned}
 \label{eq:answer3}
\end{equation}
and its inverse, as $y\to\infty$, begins $\sqrt y-\tfrac1{\sqrt2}y^{(2-\sqrt2)/2}+\tfrac{3-\sqrt2}{4}y^{(3-2\sqrt2)/2}+\OO(y^{(4-3\sqrt2)/2})$.

\subsection{What is common to all of this}

All these expansions are finite or infinite sums of \emph{power--log blocks} $w^\beta P(\log w)$, where $w\to0^+$ is a local variable ($w=y$, $w=1/x$, \dots), $\beta$ is a real exponent, and $P$ is a polynomial. Their supports lie in translates of semigroups generated by finitely many positive gaps. These semigroups are locally finite, even when the groups generated by the same gaps are dense in $\R$, so every truncation is a finite computation. The logarithmic polynomials are handled by the Euler identity $\frac{\dd}{\dd w}\bigl[w^\beta P(\log w)\bigr]=w^{\beta-1}(\beta+\tfrac{\dd}{\dd L})P(L)$; and remainders must keep their logarithmic factor, because $y^3(2\log^2y+5\log y+3)$ is not $\OO(y^3)$.

This article develops the algebra and analysis needed to construct and justify these expansions. Its starting points are the questions \cite{q1,q2} and the phase-coordinate approach of \cite{proveit}. The treatment gives coefficient formulas, branch theorems, convergence proofs and remainder estimates, then extends the construction to finite endpoints from either side, infinity, observables of an inverse, and scales involving iterated logarithms or exponentially small terms.

\subsection{Summary of results}

\begin{itemize}[leftmargin=1.6em]
\item \textbf{Scale} (Section~\ref{sec:scale}). Growth classes of power--log blocks are totally ordered, the finitely generated exponent semigroups are locally finite, and arithmetic in the truncated nonnegative-exponent algebra is exact and terminates. Real powers, logarithms and exponentials of units reduce to finite substitutions. Shifted jets require the precision rules of Section~\ref{sec:evaluator}. The common composition recurrence is $Q_{k+1}=\bigl((a-k)Q_k+Q_k'\bigr)/(k+1)$ for $(1+t)^aP(L+\log(1+t))$.
\item \textbf{Forward expansions} (Section~\ref{sec:forward}). The closure theorem gives sufficient leading-term and branch conditions for arithmetic, real powers, logarithms, exponentials and analytic composition to preserve convergent power--log expansions. Structural error propagation supplies a rigorous remainder class, which answers the second question. A positive leading logarithmic polynomial alone does not meet the constant-coefficient hypothesis for arbitrary powers and logarithms.
\item \textbf{Branch} (Section~\ref{sec:branch}). For $f(x)=ax^p(1+\sum_ix^{\delta_i}B_i(\log x))$ with $a,p\ne0$, $\delta_i>0$, the real branch is the one with $x>0$ and $x/(y/a)^{1/p}\to1$; it is unique, and the complex germ \eqref{eq:complex-germ} is explained.
\item \textbf{Lagrange--B\"urmann with a marker} (Section~\ref{sec:lagrange}). A residue proof of $\Phi(u)=\Phi(q)+\sum_N\frac{(-\eta)^N}{N!}\partial_q^{N-1}[\Phi'(q)h(q)^N]$ at a positive point, valid analytically and formally, with the transport of observables and the general-core form.
\item \textbf{Coefficients} (Section~\ref{sec:coefficients}). For every multi-index $\bk$ of weight $\sigma$, the block of $g(y)^r$ at $z^{r+\sigma}$, $z=(y/a)^{1/p}$, is
 $\frac{(-1)^{|\bk|}r}{p^{|\bk|}\bk!}\prod_{j=1}^{|\bk|-1}(\D+r+\sigma+pj)\prod_iB_i(L)^{k_i}$, $\D=\dd/\dd L$; equal weights are added. Degrees, leading logarithms (Catalan numbers), closed forms for constant coefficients (Fuss--Catalan), $\log g$, and the formal uniqueness through the triangular equation, Picard iteration and Newton iteration.
\item \textbf{Convergence} (Section~\ref{sec:convergence}). The expansion converges absolutely for small positive uniformizer $z$, including when the target tends to infinity, after substituting finitely many small parameters $z^{\delta_i}(\log z)^j$. Its coefficient bounds are $|C_{\bk}(L)|\le MK^{|\bk|}(1+|L|)^{D(\bk)}$; an explicit conditional tail majorant follows from Rouch\'e's theorem.
\item \textbf{Remainders} (Section~\ref{sec:remainders}). Truncation below any exponent cutoff leaves $z^{1+\sigma_*}C_{\sigma_*}(L)+\oo(z^{1+\sigma_*})$ where $\sigma_*$ is the first possible omitted weight and $C_{\sigma_*}$ the complete coefficient there, computed from the one-step boundary of the retained multi-index set. If that coefficient cancels, a sharper order requires further layers. Depth truncation has its own bound.
\item \textbf{Transport} (Section~\ref{sec:transport}). Forward remainders and residuals are transported to the inverse with the correct scaling $z^{1-p}$; a residual bracket gives existence and a two-sided error bound; and a theorem with a complete proof shows that finite smoothness of the perturbation, with derivative bounds in the power--log scale, suffices for the marker formula with an explicit remainder, which covers perturbations such as $x^2\sin(\log x)$.
\item \textbf{Endpoints} (Section~\ref{sec:endpoints}). Finite points from either side and $\pm\infty$ reduce to the same local inversion problem through a local variable and a signed leading power $p\ne0$; the limit $y_0$ may be finite or infinite; powers $(x-x_0)^r$ and $x^r$ of the inverse are obtained directly. Worked examples include the Wright omega function and the inverse of \eqref{eq:answer3}.
\item \textbf{Lambert coordinates and boundaries} (Section~\ref{sec:boundaries}). Leading logarithmic and exponential cores reduce to a convergent Lambert expansion in a reciprocal-logarithm coordinate, with a polynomial in its logarithm at every order. This covers $x\log x$ near $0$ and $xe^x$ at infinity. Zero gaps, flat perturbations and oscillatory coefficients require further scale choices; symbolic exponents can use depth truncation when their branch assumptions are proved.
\item \textbf{Coordinates and extended scales}. Exact changes of source and target coordinate transport inverse expansions and their errors. Weighted-support algorithms organize coefficient calculation; interval residual bounds turn local approximations into enclosures of unique roots. Further sections develop exact-core perturbations, iterated logarithmic scales, Fourier coefficients and exponentially small sectors.
\item \textbf{Special functions} (Sections~\ref{sec:special-function-expansions}--\ref{sec:special-function-adapters} and \ref{sec:exponential-normalization}--\ref{sec:barnes-inverse}). Integral and sum representations supply finite-order tail bounds; absolute envelopes retain meaning through oscillatory zeros. Logarithmic normalization gives relative Gamma and Barnes expansions. Their increasing inverses have complete polynomial reciprocal-logarithmic blocks around exact Lambert balances, with finite-order error transport and first omitted constants. These inverse results do not assert convergence as the truncation order tends to infinity.
\end{itemize}

\subsection{Reading paths and kinds of conclusion}

Sections~\ref{sec:scale}--\ref{sec:examples} develop the basic positive-gap
calculus in order: algebra, branch selection, coefficients, convergence,
remainders and examples. For a transformed inverse, continue with
Sections~\ref{sec:coordinates} and \ref{sec:source-chart-implementation};
Section~\ref{sec:inverse-function-expressions} explains how source, target
and parameter domains enter an implicit equation. Sections
\ref{sec:incremental}--\ref{sec:series-calculus} describe coefficient
refinement and precision propagation.

The extensions can then be read by scale. Logarithmic hierarchies and their
calculus appear in Sections~\ref{sec:logarithmic-scales} and
\ref{sec:reciprocal-log-operations}. Retained inverse cores are treated in
Sections~\ref{sec:core-perturbations} and
\ref{sec:exponential-core-perturbation}; finite flat sectors and their inner
precision in Sections~\ref{sec:flat-sectors} and
\ref{sec:flat-sector-operations}; Fourier coefficients in
Section~\ref{sec:fourier-coefficients}. The special-function chapters apply
these constructions to concrete functions, with each forward tail justified
separately.

The following distinctions remain in force in every scale. An expression
may satisfy more than one row, but each conclusion needs its own proof.

\begingroup
\small
\renewcommand{\arraystretch}{1.12}
\par\medskip\noindent
\begin{tabular}{@{}>{\raggedright\arraybackslash}p{0.25\linewidth}p{0.71\linewidth}@{}}
\toprule
Conclusion & What it establishes \\
\midrule
Formal identity & Coefficients agree in a stated filtered algebra. The
cutoff records discarded weights, without supplying an analytic bound. \\
Finite-order asymptotics & A specified function and branch satisfy a bound
for each fixed truncation order. The bound includes the coordinate,
logarithmic degree or other error envelope, and parameter hypotheses. \\
Convergent expansion & The infinite series converges to the selected
function on the stated domain. This is stronger than the existence of
arbitrarily long finite-order asymptotic expansions. \\
Uniform estimate & Common constants and a common domain control the
remainder while parameters vary over a specified set. Fixed-parameter
estimates alone do not allow a simultaneous parameter limit. \\
Numerical enclosure & Explicit residual and derivative bounds enclose the
selected real root at a specified target
(Section~\ref{sec:executable-certificates}). \\
\bottomrule
\end{tabular}
\par\medskip
\endgroup

\subsection{Conventions}

$\N=\{0,1,2,\dots\}$. For a local variable $w\to0^+$ we write
$L=\log w$ and $\M(w)=1+|\log w|$. Later sections explicitly introduce
other logarithmic coordinates, such as $L=\log X$ for a large retained
core $X$; their definitions apply within those sections. Unless a different
limit is stated, $\OO$ and $\oo$ refer to $w\to0^+$ with all other data
fixed. Constants and endpoint thresholds may depend on these data and on
the fixed truncation order. No uniformity in parameters is asserted unless
stated.

$\D=\dd/\dd L$ acts on polynomials in $L$. For a multi-index
$\bk=(k_1,\dots,k_m)\in\N^m$ we write $|\bk|=\sum k_i$,
$\bk!=\prod k_i!$, and $\bk\cdot\bdelta=\sum k_i\delta_i$.
Real powers of positive quantities are $w^\beta=e^{\beta\log w}$.
Real observables use the selected real branches; auxiliary complex powers
may occur in analytic proofs or Fourier coefficient representations.

Power and weight cutoffs are \emph{exclusive}: a cutoff $q$ retains
complete blocks of exponent $<q$. A depth counts a different index:
marker depth $N$ retains total marker degree at most $N$, while exponential
sector depth $N$ retains exponential degrees through $N$. Equal weights
are collected before zero blocks are removed. A finite number of vanishing
coefficients does not prove exact termination. When $P=+\infty$ is used
to mark an exact finite expression, it means zero remainder; it is not
ordinary big-$O$ notation for an error smaller than every power.
